% =============================================================
%  Springer LNCS conference draft
%  Does Explicit Devanagari G2P Still Matter for Neural TTS?
%
%  Kaustubh Pandey (2025PA17326), M.Tech CSE (AI), NSUT
%  Supervisor: Dr. Shobha Bhatt
%
%  Compile with XeLaTeX. Needs llncs.cls and splncs04.bst from
%  Springer's author kit, plus the Devanagari font if the examples
%  in Section 3.2 are kept.
%
%  Written in the first person to match the progress report. If the
%  supervisor is a co-author, change every "I" to "we" and every
%  "my" to "our" before submission.
% =============================================================
\documentclass[runningheads]{llncs}

\usepackage{iftex}
\ifPDFTeX
  \usepackage[T1]{fontenc}
  \usepackage[utf8]{inputenc}
  \newcommand{\dev}[1]{\textit{#1}}
\else
  \usepackage{fontspec}
  \newfontfamily\devfont{Noto Sans Devanagari}[Path=fonts/,
    Extension=.ttf, UprightFont=NotoSansDevanagari-Regular]
  \newcommand{\dev}[1]{{\devfont #1}}
\fi

\usepackage{graphicx}
\usepackage{booktabs}
\usepackage{amsmath}
\usepackage{url}

\begin{document}

\title{Does Explicit Devanagari Grapheme-to-Phoneme Conversion
Still Matter for Neural Text-to-Speech?}

\titlerunning{Does Explicit Devanagari G2P Still Matter?}

\author{Kaustubh Pandey\inst{1}\orcidID{0000-0000-0000-0000}}
\authorrunning{K. Pandey}

\institute{Netaji Subhas University of Technology, New Delhi, India\\
\email{kaustubh.pandey.2025pa17326@nsut.ac.in}}

\maketitle

\begin{abstract}
Devanagari does not write the inherent vowel it implies, so a phonemic front
end for Hindi has to decide where that vowel is pronounced and where it is
deleted. Writing such a front end is a month of work. Whether current neural
architectures make it unnecessary has not been measured under controlled
conditions, because published comparisons vary architecture, data, sampling
rate, vocoder and front end at the same time. I report an in-progress
testbed that varies only the input representation. Architecture, speaker,
corpus, training budget, vocoder and evaluation protocol are held fixed by
construction, with the budget enforced in mel frames so that no architecture
can be given more audio than another, and with configuration hashes that make
an unintended difference between two arms detectable rather than invisible. The
design covers FastSpeech~2 and VITS on Hindi, a nested data ladder from nine
hours to ten minutes, and Marathi as a same-script control in which medial
schwa deletion does not apply. Twenty-two runs of 100{,}000 steps are planned,
eight have completed and two are in training.

Two results are already stable. First, mel-cepstral distortion cannot resolve
this contrast at the quality level the budget affords: the two Hindi VITS arms
sit 54.3\% and 55.4\% of the distance to unrelated audio and differ by 0.89\%
of that distance, which is 0.31 standard deviations. I therefore introduce two
measures that can. One regresses excess utterance duration on the number of
rule-predicted deletion sites, giving milliseconds of excess per site, which
needs no vocoder and cannot saturate. The other partitions character error
rate from two independent connectionist temporal classification recognisers by
whether the reference word carries a rule-predicted deletion, and reports it
against the recogniser's own error rate on the real recordings. On the first
ten test utterances, both recognisers agree that errors concentrate at
\emph{medial} deletion sites for the graphemic arm and not for the phonemic
arm, by 0.1147 and 0.1150 after floor correction, while at word-final sites
the graphemic arm is marginally better. That split is what one would predict:
final deletion is inferable from the orthography, medial deletion is not.
Second, I document two measurement defects that silently invert conclusions,
one in the training objective and one in the evaluation text pipeline, and
give the diagnostic that exposes each.

\keywords{Text-to-speech \and Grapheme-to-phoneme \and Hindi \and Schwa
deletion \and Evaluation methodology \and Low-resource speech}
\end{abstract}

\section{Introduction}

Devanagari underspecifies the vowel. A consonant letter carries an inherent
\dev{अ}, and whether that vowel is pronounced depends on position, syllable
weight and morphology. A phonemic front end must therefore decide where the
schwa survives, and it will sometimes be wrong. Narasimhan et~al.
\cite{narasimhan2004schwa} gave the classical rule for Hindi, deleting a
medial schwa between two consonants that are themselves flanked by vowels;
Arora et~al. \cite{arora2020schwa} later recast the problem as supervised
sequence labelling.

Synthesis has since moved past the pipeline those rules were built for.
FastSpeech~2 \cite{ren2021fastspeech2} predicts duration, pitch and energy and
generates a spectrogram in parallel. VITS \cite{kim2021vits} maps text to
waveform through a variational autoencoder trained adversarially and learns
its own alignment. Matcha-TTS \cite{mehta2024matcha} reaches comparable
quality in fewer steps using conditional flow matching. Where a spectrogram is
produced, HiFi-GAN \cite{kong2020hifigan} turns it into audio. All three
acoustic models can be trained on characters directly, and for Indic languages
that is now common practice.

So the practical question is whether to write the front end at all. If these
architectures learn schwa deletion from data, the work is wasted and nobody
should repeat it. If they do not, the field should stop defaulting to
character input for Devanagari. Neither answer has been published, because the
comparison that would settle it has not been run with everything else held
fixed.

My contribution is that controlled comparison, plus the evaluation machinery
it turned out to need. This paper reports the design, the ten runs that have
finished, the preliminary intelligibility result, and two measurement defects
that I think are of wider interest than my own numbers, because each one
produced a plausible answer that was wrong.

\section{Related Work}

\textbf{Architectures.} FastSpeech~2 \cite{ren2021fastspeech2} depends on
durations from somewhere, traditionally an external aligner. VITS
\cite{kim2021vits} collapses the pipeline into one model and aligns internally
by monotonic alignment search, inherited from Glow-TTS \cite{kim2020glowtts}.
Matcha-TTS \cite{mehta2024matcha} is the smallest of the three and is
interesting at the low-data end of a ladder. StyleTTS~2
\cite{li2023styletts2} reports stronger naturalness than any of them and is
excluded here for compute rather than for quality.

\textbf{Alignment.} The Montreal Forced Aligner \cite{mcauliffe2017mfa} has
been the conventional duration source for the FastSpeech family. Badlani
et~al. \cite{badlani2022alignment} showed that an alignment module trained
jointly with the model, without external supervision, matches or beats
forced-aligner durations. That result decided my design, for a reason
specific to this study: an external aligner trained on a different corpus and
phone set would carry an unmeasured mismatch into exactly the arm under test,
and an aligner trained from my own dictionary would make the durations a
function of the schwa rule, which is the variable being manipulated.

\textbf{Indic resources.} The consortium recordings \cite{baby2016resources}
sit behind most Indian-language work, including both corpora used here. Kumar
et~al. \cite{kumar2023nextbillion} obtained usable quality from modest
per-language data. IndicVoices-R \cite{sankar2024indicvoicesr} scaled that for
synthesis. MMS \cite{pratap2024mms} supplies pretrained VITS checkpoints for
both of my languages and, separately, the recogniser I use as a cross-check.
These papers establish that Indic TTS works and roughly how well. None of them
separates the contribution of the front end.

\textbf{Evaluation.} Mel-cepstral distortion \cite{kubichek1993mcd} is the
standard spectral measure, and one that is partly a loudness measure if the
zeroth coefficient is retained. Pitch error on a linear scale is dominated by
the higher-pitched speaker, and pitch itself is usually taken from WORLD
\cite{morise2016world}, where this testbed uses a YIN tracker gated on frame
energy and records that choice rather than defending it as optimal. Intelligibility is increasingly assessed by
transcribing synthesised audio, often with Whisper
\cite{radford2023whisper}, and automatic predictors such as UTMOS
\cite{saeki2022utmos} and NISQA \cite{mittag2021nisqa} stand in for listening
tests. Both predictors were trained largely on English and Japanese. Section
\ref{sec:intel} argues that a sequence-to-sequence recogniser is the wrong
instrument for a schwa contrast whatever its accuracy, and that the
recogniser's own error rate on natural audio has to be reported with any
transcription-based number.

\section{Testbed}

\subsection{Corpora and Standardisation}

I use the Hindi and Marathi corpora from the SPRING Lab IndicTTS collection,
one studio speaker per language. No dataset card states hours, speaker count
or sampling rate, so I measured all three; an estimate from a 50-utterance
sample proved 19\% low for Hindi. Of 22{,}764 raw utterances, 22{,}088 survive
standardisation.

Every utterance is downmixed to mono, silence-trimmed at 40~dB below peak with
50~ms of padding retained, normalised to $-23$~LUFS \cite{itu2015bs1770},
resampled to a 22.05~kHz master and separately to 16~kHz, and peak-guarded at
0.99. Two details are deliberate. Trimming happens at the 48~kHz source, so
both outputs share their boundaries to the sample. And the 16~kHz copy is
derived from the same gained source rather than from the master, because
cascading two resamples would put a measurably different signal in front of
one architecture and that difference would later read as an architectural
effect.

Splits are frozen by salted SHA-256 and recorded in a lock file. They rebuild
byte-identically on a different machine, operating system and filesystem, which
is checked rather than assumed. The ladder is nested: the ten-minute rung is a
subset of the thirty-minute rung, which is a subset of the one-hour rung, and
so on up through five and nine hours. Independent samples at each rung would
confound data quantity with data identity.

\subsection{Front End}

The front end has four separable stages, with schwa deletion as a stage that
can be switched off independently of the others. Over 11{,}044 transcripts and
870{,}547 phone tokens it emits no unknown symbol.

Accuracy is 39 of 47 words, 83.0\%, against gold forms built from elicited
native-speaker yes/no judgements rather than from my own answer key. Four
qualifications travel with that number. Three speakers' responses were
averaged into one sheet before the merge, so inter-speaker agreement cannot be
recovered and the apparent zero variation is an artefact of the averaging.
Seven of the eight errors are speakers retaining a schwa the rule deletes,
which is citation-form reading, so 83.0\% is a lower bound under conditions
that favour retention. Four of the eight fall in one environment, word-final
schwa after a consonant cluster in Sanskrit-derived words; the configuration
has a blocking list for exactly this and it is left empty, because filling it
from these four words would fit the rule to its own test set. One elicited
answer is impossible and is left in the denominator rather than quietly
dropped.

\subsection{Holding the Budget Fixed}
\label{sec:budget}

Batching is by mel frames, charged as $|B| \times \max_i T_i$ so that padding
is paid for rather than hidden, with a cap of 12{,}000 frames per step and
100{,}000 steps for every run in the matrix. At hop 256 a step therefore
consumes exactly $12{,}000 \times 256 = 3{,}072{,}000$ samples, which is
139.32~s of audio at 22.05~kHz and 192.00~s at 16~kHz. The asymmetry is
recorded rather than corrected, because correcting it would make the step
counts differ instead.

One consequence of a fixed step budget on a nested ladder is worth stating
explicitly, since it governs how the ladder must be read. The number of passes
over the data is inversely proportional to the rung size: 430 epochs at nine
hours, 774 at five, 3{,}870 at one, 7{,}740 at thirty minutes and 23{,}220 at
ten. The lowest rungs are not undertrained. They are memorised.

Every run is described by a configuration whose hash covers all fields except
a small cosmetic set. Two arms that should differ only in input representation
therefore have hashes that differ in exactly that field, and an accidental
divergence is detectable. A separate assertion refuses to start a run whose
configuration declares a warm start that its adapter does not read; this was
added after eight runs trained from scratch while their configurations claimed
otherwise, which nothing caught until the audio was listened to.

\subsection{Matrix and Progress}

Twenty-two runs: the two Hindi arms for each of FastSpeech~2 and VITS, a
nested ladder of four rungs for each architecture, both arms of the Marathi
control, two vocoders, and four seed-repeat runs. The seed runs matter more
than their number suggests. Until the spread between two runs that differ only
in random seed is known, no gap between two arms can be reported as a finding,
and the rest of the matrix is a set of single samples with nothing to measure
them against.

Eight runs of 100{,}000 steps have completed: both Hindi arms of FastSpeech~2
and of VITS, the five-hour and one-hour FastSpeech~2 rungs, and both Marathi
FastSpeech~2 arms. The VITS seed pair is in training, which leaves twelve runs
not yet started.
FastSpeech~2 has 37.99~M parameters and VITS 83.05~M, both measured rather
than quoted.

\section{Evaluation}

\subsection{Why Mel-Cepstral Distortion Is Not Enough}
\label{sec:mcd}

I compute mel-cepstral distortion with one dynamic-time-warping alignment
shared by the spectral and pitch measures, and I report it against a chance
level obtained by rotating the reference set by one utterance, which gives the
score of mismatched audio on the same material and the same alignment.

Against that baseline the two Hindi VITS arms sit at 54.3\% and 55.4\% of the
distance to unrelated audio. They differ by 0.89\% of chance, which is 0.31
standard deviations. At the quality the budget affords, the distance of both
systems from the reference swamps the difference between them. A measure that
cannot separate the arms is not evidence that the arms do not differ, and
reporting it as such would be the error this section exists to avoid.

\subsection{Schwa Deletion in Milliseconds}
\label{sec:duration}

A retained schwa takes time. Under the rule, 32.0\% of Hindi words carry a
deletion site; the phonemic arm is handed the deletion and the graphemic arm
has to infer it. So I fit
\begin{equation}
t_{\text{syn}} - t_{\text{ref}} = \beta_0 + \beta_1 n_{\text{sites}}
\end{equation}
over the test set, where $n_{\text{sites}}$ is the number of rule-predicted
deletion sites in the input text and $\beta_1$ is milliseconds of excess
duration per site. The intercept is fitted rather than assumed so that a
uniformly slow model does not read as one retaining schwas. Deletion-site
counts come from the text, so both arms receive the same count for the same
utterance, which is the point of the design.

On synthetic data where the truth is known, a model retaining every schwa
gives $\beta_1 = 67.8$~ms at $t = 45.1$ and $r = 0.99$, and a model deleting
correctly gives $0.5$~ms at $t = 0.4$ and $r = 0.05$. The two fitted slopes
are separated by more than five standard errors at realistic noise, against
$t = 0.31$ for mel-cepstral distortion on the same contrast. The fit refuses
to return a value in the two cases where a regression would otherwise invent
one: fewer than three usable points, and no variation in site count across
utterances, where no slope is identifiable.

This measure needs no vocoder, because a mel frame count is a duration. It
therefore applies to the FastSpeech~2 arms before their vocoder has trained,
which is otherwise a hard block on every waveform metric.

\subsection{Intelligibility, and Why the Recogniser Must Be Chosen Carefully}
\label{sec:intel}

Transcribing synthesised speech and comparing the transcript with the input
text is the conventional intelligibility proxy. It is only sound if the
recogniser does not repair the errors being counted, and for this contrast
that condition excludes the strongest available recognisers. A
sequence-to-sequence decoder carries an implicit language model. Given a
slurred Hindi word it will emit the correct spelling, because the spelling is
what its training data makes likely, and a deleted or retained schwa is
precisely the kind of error such a decoder repairs. A weaker connectionist
temporal classification model, decoded greedily with no external language
model, is frame-local and reports what it heard. For this measurement it is
the better instrument.

I therefore use two CTC recognisers and never average them. IndicConformer
\cite{ai4bharat2025indicconformer} is primary: MIT licensed, built for 22
Indian languages, Devanagari output. MMS \cite{pratap2024mms} with the Hindi
adapter is the cross-check, a different architecture trained on different
data. Both are pinned to a full commit hash, which for IndicConformer fixes
the executed code and not only the weights. Disagreement between them is
reported as a result, because it is evidence about the instrument.

Character error rate leads rather than word error rate. Devanagari word
boundaries interact with the sandhi and schwa phenomena under test, so a rate
counting word boundaries partly measures my own tokenisation; one missed
schwa in a three-word sentence costs a third of the word error rate and a
tenth of the character error rate. Rates are pooled as total edits over total
reference characters rather than averaged per utterance, because short words
are where a deleted final schwa is proportionally largest.

Two further choices follow from Section \ref{sec:mcd}. The errors are
partitioned by whether the reference word carries a rule-predicted deletion,
with word-final and medial sites kept in separate classes, since the duration
measure of Section \ref{sec:duration} cannot distinguish them. And the real
recordings are transcribed too, giving the recogniser's own error rate as a
floor. Without that floor an absolute rate on synthesis has no denominator,
and the same partition on natural audio says whether the recogniser itself
struggles at deletion sites, which would otherwise be read as a property of
the models.

\section{Preliminary Results}

\subsection{Intelligibility on the Hindi VITS Arms}

Table \ref{tab:intel} gives the first ten utterances of the Hindi test split,
scored on the two VITS arms by both recognisers, with the floor above the runs.
The class denominators are small and are stated with the rates rather than left
to be inferred: 327 reference characters in 82 word tokens carry a word-final
site, 96 characters in 21 tokens carry a medial site, and 389 characters in 124
tokens carry neither, for 812 characters in total.

\begin{table}[t]
\centering
\caption{Character error rate, pooled, on the first ten Hindi test utterances.
\textsc{gt} is the recogniser's own rate on the real recordings. Classes are
by whether the reference word carries a rule-predicted deletion site. Excess is
the site class minus the no-site class. Counts in parentheses are edits over
reference characters.}
\label{tab:intel}
\begin{tabular}{llrrrr}
\toprule
Recogniser & Condition & All & Final site & Medial site & No site \\
\midrule
IndicConformer & \textsc{gt} & 0.0259 & 0.1040 & 0.0104 & 0.0411 \\
               &             & (21/812) & (34/327) & (1/96) & (16/389) \\
               & phonemic    & 0.1047 & 0.2813 & 0.0417 & 0.1208 \\
               &             & (85/812) & (92/327) & (4/96) & (47/389) \\
               & graphemic   & 0.1330 & 0.2569 & 0.1667 & 0.1311 \\
               &             & (108/812) & (84/327) & (16/96) & (51/389) \\
\midrule
MMS-1B-all     & \textsc{gt} & 0.0333 & 0.1346 & 0.0521 & 0.0540 \\
               &             & (27/812) & (44/327) & (5/96) & (21/389) \\
               & phonemic    & 0.1983 & 0.4526 & 0.2188 & 0.2262 \\
               &             & (161/812) & (148/327) & (21/96) & (88/389) \\
               & graphemic   & 0.2266 & 0.4587 & 0.3750 & 0.2674 \\
               &             & (184/812) & (150/327) & (36/96) & (104/389) \\
\bottomrule
\end{tabular}
\end{table}

Three things in that table are worth separating.

\textbf{The direction of the overall rate agrees across recognisers.} The
graphemic arm is worse by 0.0283 on both, which is 23 additional edits out of
812 in each case. The identical figure is a coincidence of small integers and
would not be expected to recur at larger $n$; the underlying counts differ
everywhere else.

\textbf{The site-specific effect is medial, not final.} Table
\ref{tab:excess} subtracts the floor from each arm's excess. At medial sites
the phonemic arm's errors avoid those words and the graphemic arm's
concentrate on them, by 0.1147 and 0.1150 for the two recognisers. At
word-final sites the graphemic arm is marginally the better of the two, by
$-0.0347$ and $-0.0351$. Two recognisers of different architecture, trained on
different data, agree on both figures to within $4 \times 10^{-4}$.

\begin{table}[t]
\centering
\caption{Excess character error rate at deletion sites, relative to words with
no site, after subtracting the recogniser's own excess on the real recordings.
Positive means errors concentrate where the rule fires.}
\label{tab:excess}
\begin{tabular}{llrrr}
\toprule
Site & Recogniser & Phonemic & Graphemic & Difference \\
\midrule
Final  & IndicConformer & $+0.0976$ & $+0.0629$ & $-0.0347$ \\
       & MMS-1B-all     & $+0.1458$ & $+0.1107$ & $-0.0351$ \\
Medial & IndicConformer & $-0.0484$ & $+0.0663$ & $+0.1147$ \\
       & MMS-1B-all     & $-0.0055$ & $+0.1095$ & $+0.1150$ \\
\bottomrule
\end{tabular}
\end{table}

That split is what the phonology predicts. Word-final schwa deletion in Hindi
is close to categorical and is inferable from the orthography alone, since a
word-final consonant simply does not carry its inherent vowel; a model trained
on characters can learn it as a statistical regularity of the script. Medial
deletion depends on syllable weight and morphology, which the graphemes
underdetermine. An explicit front end should pay where the orthography stops
being sufficient, and on this evidence that is where it pays.

\textbf{The nuisance variable sits in the other class.} On the real
recordings both recognisers err substantially more per character on
word-final-site words than on words with no site: 34/327 against 16/389, and
44/327 against 21/389. Those words are longer, rarer and carry more consonant
clusters, so part of the apparent final-site concentration has nothing to do
with synthesis. The floor's \emph{medial} excess, by contrast, is $-0.0307$ and
$-0.0019$, which is to say nothing. The confound is concentrated in the class
where the effect is absent and absent from the class where the effect is
found.

\textbf{What this is not.} The medial class is 96 characters in 21 word
tokens, and the counts behind the headline ratio are 16 edits against 4 for
IndicConformer and 36 against 21 for MMS. A single utterance could be carrying
most of it. No confidence interval is computable from one sample of ten
utterances, and none is claimed. The result is reported here because the
agreement between two independent recognisers is itself informative, and
because the measure's behaviour is worth describing before the full split is
scored.

\section{Two Measurement Defects}

Both of the following produced a plausible number that was wrong in a
direction that mattered. I give each with the diagnostic that exposed it,
because neither would have been caught by a test of the code's correctness.

\subsection{A Training Loss That Is Mostly a Pitch Error}
\label{sec:loss}

The FastSpeech~2 runs end at a final training loss of 314.60 at nine hours and
299.85 at five hours, and then 10.04 at one hour. A factor of 29.9 at the
one-hour rung has no reading as overfitting, because an overfitting curve has
no cliff in it, and the ladder could not be plotted against training loss
until the number was explained.

The objective is a sum of five terms, and two of them are mean squared errors
in physical units: fundamental frequency in hertz and the Euclidean norm of
each linear spectrogram frame, both weighted at 0.1. Neither is normalised
anywhere in my pipeline, because the upstream implementation normalises them
inside a dataset class that this testbed does not use. Measured on twelve
utterances of each rung, through the same feature code the models trained on,
a mean predictor's share of the loss is as shown in Table \ref{tab:loss}.

\begin{table}[t]
\centering
\caption{Contribution of each loss term at a mean predictor, 22.05~kHz, twelve
utterances per column. The mel term is the only one that measures spectral
quality.}
\label{tab:loss}
\begin{tabular}{lrr}
\toprule
Term & One-hour rung & Nine-hour only \\
\midrule
$1.0 \times$ mel $L_1$ against a mean predictor & 1.814 & 1.972 \\
$0.1 \times \operatorname{var}(f_0$ in Hz$)$ & 569.131 & 653.995 \\
$0.1 \times \operatorname{var}($frame energy$)$ & 45.381 & 45.647 \\
\midrule
Sum & 616.326 & 701.614 \\
\bottomrule
\end{tabular}
\end{table}

So a mean predictor scores roughly 617 to 702, and the mel term is under one
per cent of it. The nine-hour run's 314.60 implies a pitch error of about
51.7~Hz against a speaker standard deviation of 80.9~Hz. The one-hour run's
10.04 implies about 10.0~Hz. Since the rungs are nested subsets of one
speaker, the units are identical at every rung; what differs is that 541
utterances seen 3{,}870 times have their fundamental frequency memorised. The
mel term is 0.6\% of the nine-hour loss and 18\% of the one-hour loss, so the
same number does not mean the same thing at two rungs.

The consequence for the ladder is that it must be plotted against held-out
metrics and not against training loss, with loss curves confined to an
appendix carrying this decomposition beside them. The consequence for
FastSpeech~2 is that its mel decoder is optimised against under one per cent of
the gradient signal while an unnormalised pitch head takes most of the rest,
which is a sufficient explanation for the poor quality of those arms and is a
property of the objective rather than of the data or the architecture.
Normalising the two terms would change the configuration hash of ten runs, so
I record the defect and keep the objective fixed for the remainder of the
matrix rather than making the arms incomparable mid-study. The per-term
breakdown is now logged at every step, so the claim is checkable from the runs
instead of reconstructed from a filter bank.

\subsection{Punctuation That Moved Words Between Classes}
\label{sec:punct}

No CTC recogniser emits punctuation, and the reference transcripts are
punctuated. On the first two ground-truth utterances of the Hindi test split,
41\% and 48\% of the measured character error was commas and a full stop that
the recogniser was never going to produce. On its own that is an inflation
affecting all conditions equally.

The second effect is not benign. The site classifier asks whether a
rule-predicted deletion falls on the \emph{last} segment of a word, and a
trailing comma is one more segment. All six word-final-site words sampled from
the test split were reclassified from final to medial by appending a comma. The
word-final class had therefore been excluding every clause-final word, which is
the position where final schwa deletion is most audible and most prosodically
marked, and the medial class had been contaminated with them.

The symptom was visible in the first table I produced and I misread it as
small-sample noise: the punctuation share of characters was 0.0000 in the
final class against 0.0690 and 0.0561 in the other two. Zero out of 283
characters is not chance. After stripping punctuation, 8 word tokens and 44
characters move from the medial class back to the final class in ten
utterances, and the floor's character error rate falls by roughly half.

The correction changed the conclusion rather than its precision. Before it, the
two recognisers disagreed about the final-site effect, giving $-0.0038$ and
$+0.1871$; after it they agree at $-0.0347$ and $-0.0351$, and the effect
relocates to the medial class. Two notes for anyone building a similar
pipeline. Strip only Unicode punctuation categories, since combining marks are
not punctuation and the nukta, anusvara and visarga are exactly the contrasts a
Devanagari study is about. And derive the class partition and the error rate
from one shared normalisation, since two copies of that chain are two
conventions.

\section{Threats to Validity}

\textbf{One speaker per language.} Every result is within-speaker. Nothing here
separates a property of Hindi from a property of this speaker's Hindi.

\textbf{The front end is 83.0\% accurate.} Both arms inherit its errors
asymmetrically: the phonemic arm is handed a wrong phoneme where the rule is
wrong, while the graphemic arm is handed the orthography and may or may not
reproduce the correct form. An effect measured against rule-predicted sites is
measured against the rule, not against the truth.

\textbf{The budget asymmetry of Section \ref{sec:budget}} gives the 16~kHz arm
192.00~s of audio per step against 139.32~s at 22.05~kHz. It is recorded, not
corrected.

\textbf{Intelligibility is an ASR proxy.} Even a CTC recogniser has acoustic
priors, and the floor reported with every rate is the available defence rather
than a complete one. A listening test is planned and not yet run.

\textbf{The seed floor is not yet measured.} Until it is, the gaps in Tables
\ref{tab:intel} and \ref{tab:excess} have no scale against them. This is the
single largest gap in the present evidence and the reason none of these
numbers is offered as a finding.

\textbf{Marathi is a control, not a replication.} It shares the script and
most of the phoneme inventory \cite{ohala1999hindi, dhongde2009marathi} and
does not apply Hindi's medial deletion rule, which is what makes it able to
separate script from phonology. It is not a second test of the Hindi result.

\section{Conclusions and Remaining Work}

The testbed holds architecture, speaker, corpus, budget, vocoder and
evaluation fixed and varies only the input representation, with the budget
enforced in mel frames and configuration hashes that make an unintended
difference detectable. Eight of twenty-two runs have finished and two are in
training.

Mel-cepstral distortion does not resolve this contrast at the quality the
budget affords, at 0.31 standard deviations, so I report two measures that
can: milliseconds of excess duration per rule-predicted deletion site, which
needs no vocoder and cannot saturate, and a site-partitioned character error
rate from two independent CTC recognisers, floor-corrected against their own
error on the real recordings. On ten utterances both recognisers agree that
the graphemic arm's errors concentrate at medial deletion sites and not at
final ones, which is the asymmetry the phonology predicts, and which no global
error rate would have shown.

The two defects of Section \ref{sec:loss} and Section \ref{sec:punct} are
offered as the more portable contribution. A training loss whose dominant
terms are in physical units is not comparable across data scales even when the
data is nested, and a site-partitioned error rate computed on punctuated
references partitions on the punctuation. Each produced a confident wrong
answer, and in the second case the wrong answer was the one the hypothesis
predicted.

What remains is the twelve outstanding runs, the seed-variance floor that
gives every gap its scale, the full test split in place of ten utterances with
a bootstrap confidence interval over utterances, the vocoder that unblocks
waveform metrics for the FastSpeech~2 arms, the Marathi control, correlation
of predicted mean opinion score against human ratings on Devanagari, a
listening test, and real-time factor on consumer hardware. The testbed, the
checkpoints and a speaker-validated contested-schwa stress-test set will be
released.

\begin{credits}
\subsubsection{\ackname} I thank Dr.~Shobha Bhatt for supervision, and the
native speakers who provided the schwa judgements of Section 3.2.

\subsubsection{\discintname} The author has no competing interests.
\end{credits}

\begin{thebibliography}{99}

\bibitem{ren2021fastspeech2}
Ren, Y., et al.: FastSpeech 2: Fast and high-quality end-to-end text to
speech. In: ICLR (2021). arXiv:2006.04558

\bibitem{kim2021vits}
Kim, J., Kong, J., Son, J.: Conditional variational autoencoder with
adversarial learning for end-to-end text-to-speech. In: ICML (2021).
arXiv:2106.06103

\bibitem{mehta2024matcha}
Mehta, S., et al.: Matcha-TTS: A fast TTS architecture with conditional flow
matching. In: ICASSP (2024). arXiv:2309.03199

\bibitem{kong2020hifigan}
Kong, J., Kim, J., Bae, J.: HiFi-GAN: Generative adversarial networks for
efficient and high fidelity speech synthesis. In: NeurIPS (2020).
arXiv:2010.05646

\bibitem{kim2020glowtts}
Kim, J., et al.: Glow-TTS: A generative flow for text-to-speech via monotonic
alignment search. In: NeurIPS (2020). arXiv:2005.11129

\bibitem{li2023styletts2}
Li, Y.A., et al.: StyleTTS 2: Towards human-level text-to-speech through style
diffusion and adversarial training with large speech language models. In:
NeurIPS (2023). arXiv:2306.07691

\bibitem{badlani2022alignment}
Badlani, R., et al.: One TTS alignment to rule them all. In: ICASSP (2022).
arXiv:2108.10447

\bibitem{mcauliffe2017mfa}
McAuliffe, M., et al.: Montreal Forced Aligner: Trainable text-speech
alignment using Kaldi. In: Interspeech, pp. 498--502 (2017)

\bibitem{pratap2024mms}
Pratap, V., et al.: Scaling speech technology to 1,000+ languages. Journal of
Machine Learning Research \textbf{25} (2024). arXiv:2305.13516

\bibitem{kumar2023nextbillion}
Kumar, G.K., et al.: Towards building text-to-speech systems for the next
billion users. In: ICASSP (2023). arXiv:2211.09536

\bibitem{sankar2024indicvoicesr}
Sankar, A., et al.: IndicVoices-R: Unlocking a massive multilingual
multi-speaker speech corpus for scaling Indian TTS. In: NeurIPS Datasets and
Benchmarks (2024). arXiv:2409.05356

\bibitem{baby2016resources}
Baby, A., Thomas, A.L., Nishanthi, N.L., et al.: Resources for Indian
languages. In: Proc. Text, Speech and Dialogue (CBBLR Workshop) (2016)

\bibitem{narasimhan2004schwa}
Narasimhan, B., Sproat, R., Kiraz, G.: Schwa-deletion in Hindi text-to-speech
synthesis. International Journal of Speech Technology \textbf{7}(4),
319--333 (2004)

\bibitem{arora2020schwa}
Arora, A., Gessler, S., Hundt, N.: Supervised grapheme-to-phoneme conversion
of orthographic schwas in Hindi and Punjabi. In: ACL (2020). arXiv:2004.10353

\bibitem{ohala1999hindi}
Ohala, M.: Hindi. In: Handbook of the International Phonetic Association, pp.
100--103. Cambridge University Press (1999)

\bibitem{dhongde2009marathi}
Dhongde, R.V., Wali, K.: Marathi. John Benjamins (2009)

\bibitem{kubichek1993mcd}
Kubichek, R.: Mel-cepstral distance measure for objective speech quality
assessment. In: IEEE Pacific Rim Conf. on Communications, Computers and Signal
Processing, pp. 125--128 (1993)

\bibitem{morise2016world}
Morise, M., Yokomori, F., Ozawa, K.: WORLD: A vocoder-based high-quality
speech synthesis system for real-time applications. IEICE Trans. on
Information and Systems \textbf{E99-D}(7), 1877--1884 (2016)

\bibitem{radford2023whisper}
Radford, A., et al.: Robust speech recognition via large-scale weak
supervision. In: ICML (2023). arXiv:2212.04356

\bibitem{saeki2022utmos}
Saeki, T., et al.: UTMOS: UTokyo-SaruLab system for VoiceMOS Challenge 2022.
In: Interspeech (2022). arXiv:2204.02152

\bibitem{mittag2021nisqa}
Mittag, G., et al.: NISQA: A deep CNN-self-attention model for
multidimensional speech quality prediction with crowdsourced datasets. In:
Interspeech (2021). arXiv:2104.09494

\bibitem{itu2015bs1770}
International Telecommunication Union: Algorithms to measure audio programme
loudness and true-peak audio level. Recommendation ITU-R BS.1770-4 (2015)

\bibitem{ai4bharat2025indicconformer}
AI4Bharat: IndicConformer 600M multilingual, automatic speech recognition for
22 Indian languages. Hugging Face model repository, revision
\texttt{e9b71b36}, MIT licence.
\url{https://huggingface.co/ai4bharat/indic-conformer-600m-multilingual}

\end{thebibliography}

\end{document}
